\setcounter{page}{7}
\begin{center}
    \vspace{0.5ex}
    \setfont{3}{1.5}\bfseries 摘要
    \vspace{0.5ex}
\end{center}

\setfont{-4}{1.5}


这篇论文基于高能效高精度全动态放大器研究项目，首先回顾了动态放大器作为低功耗放大器具有
的全动态功耗、高动态性能的优势，分析了传统固定偏置式电平移位方法在低电压环境下增益不足、线性下降、共模漂移敏感等问题的根源。
针对这些问题，本文提出了一种基于低压浮动反相放大器（LVFIA）和相关电平移位
（Correlated Level Shifting, CLS）技术的创新架构设计方法，
通过理论分析、仿真验证和实验测试，系统评估了CLS在提升动态放大器性能方面的有效性，
并与传统固定偏置式电平移位方法进行了对比。结果表明，CLS能够显著提升动态放大器的增益和线性度，
同时降低共模漂移敏感性，为未来低电压模拟电路设计提供了新的思路和方法。

\par
\noindent{\bfseries 关键词：}动态放大器；相关电平移位；低压模拟集成电路设计

\newpage
\setcounter{page}{9}

\begin{center}
    \vspace{0.5ex}
    \setfont{3}{1}\bfseries Abstract
    \vspace{0.5ex}
\end{center}

Based on the research project of high-efficiency and high-precision full-dynamic amplifier, 
this paper first reviews the advantages of dynamic amplifiers as low-power amplifiers, 
including full-dynamic power consumption and high dynamic performance. 
It analyzes the root causes of issues such as insufficient gain, 
linearity degradation, and common-mode drift sensitivity in 
traditional fixed-bias level shifting methods under low voltage conditions.
Addressing these issues, this paper proposes an innovative architectural 
design method based on Low Voltage Floating Inverting Amplifier (LVFIA) and 
Correlated Level Shifting (CLS) technology. Through theoretical analysis, 
simulation verification, and experimental testing, 
the effectiveness of CLS in enhancing the performance of dynamic 
amplifiers is systematically evaluated, 
and it is compared with traditional fixed-bias level shifting methods. 
The results show that CLS can significantly improve 
the gain and linearity of dynamic amplifiers while reducing common-mode drift sensitivity, 
providing new ideas and methods for future low-voltage analog circuit design.

\noindent{\bfseries Keywords: }Dynamic Amplifier；Correlated Level Shifting；Low-Voltage Analog IC Design